\section{Synthetic Network Topology and Flow-to-Node Mapping Sensitivity (Experiments E10 and E18)}
\label{app:topology_sensitivity}

Because raw CICIoT2023 flow traces do not include physical network topology annotations or router graph matrices, topological centrality metrics are modeled via a synthetic enterprise scale-free network topology generated according to the Barabási-Albert model (Experiments E10 and E18). The synthetic network comprises $N = 120$ nodes with attachment parameter $m = 2$ directed edges per newly introduced node, producing 236 directed communication edges and average node degree 3.93.

Personalized PageRank centrality is computed with teleportation parameter $\alpha = 0.85$ via power iteration to tolerance $10^{-6}$. Centrality scores across the synthetic topology span from minimum $R_{\min} = 0.3853$ to maximum $R_{\max} = 0.9000$, producing empirical mean $\mu_R = 0.4189$ and standard deviation $\sigma_R = 0.0721$ (median 0.3905, 95th percentile 0.5014). Exactly two nodes exceed the critical backbone hub threshold $R \ge 0.75$: Node 0 ($R = 0.9000$, degree 41) and Node 4 ($R = 0.8089$, degree 28), yielding individual node criticality $Z = 4.59\sigma$ ($p = 2.2 \times 10^{-6}$) on the empirical graph distribution (theoretical target $Z = 4.125\sigma$ under $\mu=0.42, \sigma=0.08$). We emphasize that computing a batch-mean $Z = 41.25\sigma$ via the standard Central Limit Theorem is an invalid artifact because topological risk $R$ is assigned per ingress node and remains identical across all 100 flows within any single monitoring batch, violating the independent-flow assumption. Ingress monitoring batches are mapped to nodes via round-robin index allocation $n(b) = b \pmod{120}$. Table~\ref{tab:l5_population} reports the operational escalation of core hub nodes into Level $L_5$ under flood saturation versus nominal traffic.

\begin{table*}[t]
\caption{Level $L_5$ (Critical Hub Containment) empirical posture population under core backbone SDN hub monitoring ($R \ge 0.75$, 1,000 batches, 100,000 flows, Experiment E8/E17).}
\label{tab:l5_population}
\centering
\small
\begin{tabular}{lcccccc}
\toprule
\textbf{Operational Traffic Scenario} & \textbf{Batches} & \textbf{Posture Level} & \textbf{Attack Flows} & \textbf{Benign Flows} & \textbf{Attack Pct} & \textbf{Mean Trust} \\
\midrule
Acute Flood Saturation (97.6\% Attack) & 1,000 & $L_5$ (Critical Hub Flood) & 97,600 & 2,400 & 97.60\% & 10.70 \\
Nominal Traffic Baseline (0\% Attack) & 1,000 & $L_1$ (Nominal Core Hub) & 0 & 100,000 & 0.00\% & 92.40 \\
\bottomrule
\end{tabular}
\end{table*}

To test structural robustness, Experiment E18 evaluated sensitivity across three base topology densities by adjusting preferential attachment to yield $\mu_R \in \{0.30, 0.42, 0.55\}$. Across all three topological regimes, posture allocations on edge streams remain strictly invariant ($L_0$: 12.83\%, $L_1$: 17.17\%, $L_2$: 9.17\%, $L_3$: 30.83\%, $L_4$: 30.00\%, $L_5$: 0.00\%), confirming that edge monitoring is resilient to underlying graph density shifts.
